\subsection{Full-character traces and low Verlinde compatibility}
\label{subsec:round3-ambient}
The full character field can contain useful information that disappears
under projection to the Frobenius--Perron dimension field.  We retain this
information by applying the ambient trace construction of
Proposition~\ref{prop:ambient-trace-domain} to a cyclotomic field containing
all characters.  The following formula makes this extension inexpensive.

\begin{lemma}[Direct cyclotomic relative trace]
\label{lem:round3-relative-trace}
Let $K=\QQ(\zeta_F)$ and let $\theta=\zeta_m^b$ have exact order $m$.
Put $d=\gcd(m,F)$ and $r=m/d$.  If $\gcd(r,F)>1$, the normalized relative
trace of $\theta$ to $K$ is zero.  Otherwise it is
\[
 \frac{\Tr_{K(\theta)/K}(\theta)}{[K(\theta):K]}
 =\frac{\mu(r)}{\varphi(r)}\,\zeta_d^{\,b r^{-1}},
\]
where the inverse is taken modulo $d$, and the last factor is $1$ when
$d=1$.
\end{lemma}
\begin{proof}
If a prime divides both $r$ and $F$, its exponent in $m$ exceeds its
exponent in $F$.  The last cyclotomic layer at this prime fixes $K$ and
the other primary components; the corresponding primitive root sum is
zero.  Transitivity of normalized trace proves the first assertion.
Otherwise $r$ is coprime to $F$, and hence to $d$.  The Chinese remainder
identity writes
\[
 \zeta_m^b=\zeta_d^{\,b r^{-1}}\zeta_r^{\,b d^{-1}}.
\]
The first factor lies in $K$.  Averaging the second factor over the
cyclotomic Galois group gives the primitive root sum $\mu(r)$ divided
by $\varphi(r)$.  Coprimality ensures that adjoining $K$ does not change
this average.
\end{proof}

Apply the lemma separately to the powers $\theta^j$.  With the finite
bound $B$ of Proposition~\ref{prop:ambient-trace-domain}, this computes
the complete joint trace domain directly in $K$, without constructing
$\QQ(\zeta_B)$.  An independent implementation compares this formula
with explicit Galois averaging in compositum fields: all $1056$ exact
comparisons, covering $176$ conductor--order pairs, agree.

\begin{proposition}[Full-character moment compatibility]
\label{prop:round3-full-character}
Suppose a commutative fusion ring has a pseudo-unitary categorification.
Let $D$ be its global dimension and let $K=\QQ(\zeta_F)$ contain its
full character table.  For every possible low induction frame, the
nonnegative counts of high restriction columns and joint twist traces
satisfy the following conditions:
\begin{enumerate}[label=\textup{(\roman*)}]
\item the induction Gram identities and the central indicator formulas;
\item for every simple $X$, every tested order $n$, and every proper
 divisor $r$ of $n$, the cyclic rotation character identity
\[
 \Tr(E_{X,n}^{r})=\nu_{n/r}(X^r),\qquad
 \dim\operatorname{Hom}(\one,X^n)=[X^n:\one];
\]
\item the low-pair modular equations and the integrality and Galois
 compatibility of the joint fusion eigenvalues;
\item equality of the multiplicities of paired types
 $(\lambda,\tau)$ and $(\sigma\lambda,\sigma^2\tau)$, where $\lambda$
 is the joint low fusion eigenvalue vector and $\tau$ the joint
 normalized twist-trace vector.
\end{enumerate}
Thus exact infeasibility of all these finite systems excludes a
pseudo-unitary categorification.
\end{proposition}
\begin{proof}
The Gram identities and the indicator formula are the established
induction constraints; the latter are applied coefficientwise in $K$
using \cite[Theorem~4.1]{NS}.  The cyclic rotation operator has $n$th
roots of unity as eigenvalues, and the divisor-power identity follows
by grouping the tensor factors into blocks of length $r$; see
\cite{NSpiv}.  Introduce nonnegative variables for these rotation
eigenvalue multiplicities and apply normalized relative trace.  The
resulting equations are rational after expansion in a power basis.

For low central objects $A,B$, both twists are one, so the center
identity $STS=DT^{-1}ST^{-1}$ gives the low-pair equations.
Their joint eigenvalues are algebraic integers and occur with
Galois-invariant multiplicities.  Under a Galois automorphism the
fusion eigenvalues transform by $\sigma$, whereas the canonical center
twists transform by $\sigma^2$.  Normalized relative trace commutes
with these automorphisms, proving the paired-type equality.  If an
ambient automorphism requires a lift to a compositum, the resulting
normalized trace is independent of that lift.  The twist-one low
objects contribute the known constant terms.  Relaxing integer
multiplicities to nonnegative real variables enlarges the set of
possible solutions, so an exact separating functional proves the
claimed necessary obstruction.
\end{proof}

There is a complementary refinement using low fusion products.  Write
$v_Z$ for the restriction column of a central simple $Z$ and
$d(v)=\sum_i v_i\FPdim(b_i)$.  For a low central simple $A_a$ associated
to a character $\chi_a$ of formal codegree $f_a$, put
\[
 t_a(v)=\frac{D}{f_a}\frac{\chi_a(v)}{d(v)}.
\]
Then $t_a(v_Z)=S_{A_a,Z}/d_Z$.

\begin{corollary}[Self-dual low-product constraints]
\label{prop:round3-low-verlinde}
For self-dual input, all low central simples are self-dual. The general
word and restriction formulas in Section~\ref{sec:unified-word-gram}
therefore apply with $a^*=a$ and real $\lambda_a(Z)=t_a(v_Z)$.
In particular the cubic moments are the low fusion coefficients, and
summing the nonnegative high coefficients gives the restriction-product
equations and their coordinatewise upper bounds.
\end{corollary}
\begin{proof}
The input character values are real. Hence the corresponding low
$S$-rows are real, and $S_{a^*,Z}=\overline{S_{a,Z}}$ together with
invertibility of $S$ forces $a^*=a$. The assertion follows from the
general formulas, retaining the same bounded high coefficient variables.
\end{proof}

The implementation may retain integer variables and test congruences.
Only rows with integer-count variables are used for a congruence
witness.  In the applications below, however, nonnegative linear
infeasibility already suffices: each stored ray $y$ satisfies
$y^TA\ge0$ and $y^Tb<0$ over the integers.  No numerical optimizer is
used by the final verifier.

\begin{theorem}\label{thm:round3-m4-211-243}
The rings $\RR{4}{211}$ and $\RR{4}{243}$ have no complex fusion
categorification.
\end{theorem}
\begin{proof}
For $\RR{4}{211}$, the simple dimensions are
$(1,1,d,2d+1,2d+1,d)$, where $d=2+\sqrt6$, and
$D=120+48\sqrt6$.  Its full character field is
$\QQ(\sqrt2,\sqrt3)$, so take $F=24$.  At height three the Cauchy
prime support $\{2,3\}$ gives $B=144$ and $53$ joint trace types.
The exhaustive low enumeration has $130$ frames.  Only two frames
have relevant character assignments: each has two assignments and
four second-indicator choices.  All $16$ systems are excluded by
exact rays from Proposition~\ref{prop:round3-full-character}.
Here even the weaker equality of total joint-type Galois
multiplicities suffices; the coupled trace refinement is not needed.

For $\RR{4}{243}$, the dimensions are $(1,1,d,d+1,d+1,d)$ with
$d=7+\sqrt{57}$ and $D=456+60\sqrt{57}$.  The full character field is
quadratic.  The height-six systems use
Proposition~\ref{prop:round3-low-verlinde}, including restriction-product
bounds.  There are $210$ low frames.  Four relevant frames contribute
respectively two, one, one, and two character assignments, each with
eight second-indicator choices.  The resulting $48$ systems all have
exact separating rays.

Finally, general scope is supplied independently by unrestricted
signed sphericalization.  For the two rings the complete signed
enumerations have respectively $128$ and $1584$ tables, in $4$ and
$53$ sign-gauge orbits.  Exactly one orbit in each case is gauge
isomorphic to the original tensor; the remaining character branches
are rejected by respectively $8$ and $176$ exact arithmetic
obstructions.  The surviving nonzero dimension characters belong to
the Frobenius--Perron Galois orbit.  Thus any complex categorification would admit a pseudo-unitary
Galois conjugate, contradicting the preceding exclusions.
\end{proof}

The full certificates and exact replay reports are retained in
\filename{research/round3/m4}.  The verifiers
\filename{software/round3_m4.py} and
\filename{software/round3_verlinde.py} reconstruct the complete branch
domains from the input tensors.  Both new necessary systems are also
applied to the entire frozen $22$-ring frontier and all $54$ positive
model controls, with explicit eligibility and time-budget outcomes;
a timeout is inconclusive.  In particular, $\RR{4}{155}$ and
$\RR{4}{629}$ survive the stronger full-character height-six systems
with coupled Galois spectrum/twist counts.  These surviving systems
are not asserted to be modular data or categorifications.
